%% Dans ce fichier, les �ventuels accents sont cod�s
%% en ISO-8859-15 (aka latin9)
%% La ligne suivante donne trois exemples 
%% e accent aigu : �, e accent grave : �, e-dans-l'o : �
%%%%%ENTREES%%%%
\begin{tabarg}
\multicolumn{3}{l}{\titretab {Entr\'ees}}\\
\hline
$X$ & Vecteur, \texttt{float[$N_{level}$]} & Signal mesur� \\
$Z_X$ & Coeff., \texttt{String} & Distances au mobile des mesures $X$ (\$Arg1.bornes) (m)\\
$N_{level}$ & Coeff., \texttt{int} & Nombre de niveaux (\$Arg1.array\_dim) \\
$D$ & Coeff., \texttt{int} & 1 si mesure vers le haut, -1 si mesure vers le bas  \\
$\phi$ & Vecteur, \texttt{float} & Roulis (degr�s)\\
$\theta$ & Vecteur, \texttt{float} & Tangage (degr�s)\\
$Alt$ & Vecteur, \texttt{float} & Altitude de la plateforme (m)\\
\end{tabarg}

\medskip
%%%%%SORTIES%%%%
\begin{tabarg}
\multicolumn{3}{l}{\titretab {Sorties}}\\
\hline
$Alt_X$ & Vecteur, \texttt{float[$N_{level}$]} & Altitude des mesures X (m)\\
\end{tabarg}
